\input{preamble.tex}

\begin{document}
\pagestyle{fancy}
\fancyhead[L]{Seconde}
\fancyhead[C]{\textbf{DS n°3b — fonctions numériques}}
\fancyhead[R]{\today}

\reversemarginpar

\null\vspace{-30pt}
Nom / Prénom : \\

Consignes particulières : 
\begin{itemize}[label=$\bullet$]
	\item 
	La calculatrice est {interdite}.
	\item 
	Toute trace de recherche est prise en compte.
%	\item 
%	L'abbréviation $\tq$ signifie ``tel que''.
\end{itemize}

\hrule

\exe{12}{
Résoudre dans $\R$ les équations suivantes :
\begin{multicols}{3}
\begin{enumerate}[label=(\alph*)]
\item $8x - 40 = 0$
\item $4x - 7 = 1$
\item $5x + 18 = 3$
\item $4x + 2 = x - 1$
\item $3x + \frac25 = \frac83$
\item $\frac{x}4 + 2 = \frac15$
\item $\frac92 x - 3 = \frac13$
\item $\frac53 x - 4 = 16$
\item $\frac73 x - 2 = \frac13$
\item $2(x + 4) = 3 \times \frac56$
\item $25(\frac{x}5 - \frac25) = -5$
\item $2x - \frac56 = \frac{x}3 + \frac12$
\end{enumerate}
\end{multicols}
}{exe:1}{

}

\exe{3}{
On considère le programme de calcul suivant, dont la dernière étape n'est pas entièrement connue.
\begin{enumerate}[label=(\roman*)]
\item On prend un nombre réel.
\item On mutiplie ce nombre par $\frac32$.
\item On ajoute .... à ce nombre
\end{enumerate}
On modélise ce programme de calcul par une fonction $f$ définie sur $\R$.
\begin{enumerate}
\item On sait que $f(2) = 2$. Compléter l'étape (iii) du programme de calcul.
\item Donner les éventuels antécédents de $-1$ par $f$.
\item Représenter une partie de la courbe représentative de la $f$ dans un repère.
\end{enumerate} 
}{exe:2}{

}

	
\exe{3}{
On considère la fonction $g$ définie par $g(x) = \frac6{x} + 3$.

\begin{enumerate}
\item Peut-on calculer $g$ pour toutes les valeurs de $x \in \R$ ? En déduire le domaine de définition de $g$.
\item Donner l'image de 4 par $g$.
\item Donner les éventuels antécédents de 0 par la fonction $g$.
\end{enumerate}
}{exe:3}{

}

\exe{2}{
On considère les intervalles $I_1 = ] -\infty ; 6]$ et $I_2 = ]-8 ; 9[$ et $x$ une variable réelle.
\begin{enumerate}
\item Traduire l'appartenance de $x$ aux intervalles $I_1$ et $I_2$ en termes d'inégalités.
\item Donner une condtion pour que $x$ appartienne à l'intervalle $I_1$ \textbf{ET} à l'intervalle $I_2$.
\end{enumerate}
}{exe:4}{

}

\exe{}{
(Exercice bonus) \\
Déterminer une fonction numérique $g$, définie sur $\R$ 
et pour laquelle l'image des entiers -2, 1, 1 et 2, est un nombre rationnel $\textbf{non décimal}$.
}{exe:b}{

}

%%%%%%%%%%%%
%
%\newpage
%\fancyhead[C]{\textbf{Solutions}}
%\shipoutAnswer

\end{document}
 